\chapter{Marco teórico}
\label{cap:marco-teorico}

Este capítulo presenta los fundamentos teóricos sobre los que se sustenta el
trabajo y revisa el estado del arte de las disciplinas implicadas. Se abordan,
en primer lugar, el middleware robótico y el problema de la localización y el
mapeado; a continuación, la navegación autónoma y la planificación de rutas de
cobertura; después, la detección de objetos mediante visión por computador; y,
por último, las arquitecturas orientadas a eventos. El capítulo concluye
situando la propuesta respecto a los trabajos relacionados.

%% =============================================================================
\section{Robótica móvil autónoma y ROS 2}
\label{sec:ros2}

\gls{ros} es un middleware de código abierto para el desarrollo de aplicaciones
robóticas que proporciona un modelo de comunicación entre procesos, herramientas
de construcción y un ecosistema amplio de paquetes reutilizables. Su primera
generación se introdujo con el objetivo de favorecer la reutilización de
software en investigación robótica \parencite{quigley2009ros}. Sin embargo, su
diseño, articulado en torno a un nodo maestro que coordinaba el descubrimiento
entre procesos, presentaba limitaciones para los entornos de producción, los
sistemas con varios robots y las aplicaciones con requisitos de tiempo real. La
segunda generación, \gls{ros}, rediseña la arquitectura para superar esas
carencias: elimina el maestro centralizado y delega el descubrimiento y el
transporte de mensajes en el estándar \gls{dds}, lo que aporta una comunicación
descentralizada, control de la calidad de servicio de los enlaces y mejor
soporte para despliegues distribuidos \parencite{macenski2022ros2}.

El modelo de programación de \gls{ros} se organiza en torno a \textit{nodos},
procesos independientes que encapsulan una funcionalidad concreta y se comunican
mediante el paso de mensajes. Los nodos intercambian flujos de datos a través de
\textit{topics} con un patrón de publicación-suscripción, invocan operaciones
puntuales mediante \textit{servicios} de tipo petición-respuesta y gestionan
tareas de larga duración, con realimentación y posibilidad de cancelación,
mediante \textit{acciones}. Sobre el transporte \gls{dds}, cada enlace puede
ajustar parámetros de calidad de servicio, como la fiabilidad de la entrega o la
persistencia de los últimos valores publicados, lo que permite adaptar la
comunicación tanto a sensores de alta frecuencia como a eventos esporádicos.

Más allá del modelo de comunicación, \gls{ros} ofrece un sistema de construcción
y empaquetado, un mecanismo de parámetros para configurar los nodos en tiempo de
ejecución y un sistema de lanzamiento que describe cómo arrancar y orquestar
conjuntos de nodos. Esta separación en componentes débilmente acoplados favorece
la modularidad, permite reutilizar paquetes mantenidos por la comunidad y hace
posible sustituir una pieza sin alterar el resto del sistema. 
El diseño de este trabajo sigue ese principio de forma explícita, donde cada 
capacidad del robot (telemetría, percepción, navegación, distribución de secretos y
patrulla) reside en un nodo especializado y encapsulado.

%% =============================================================================
\section{Localización y mapeado simultáneos (SLAM)}
\label{sec:slam}

Para navegar de forma autónoma, un robot necesita conocer tanto el mapa del
entorno como su propia posición dentro de él. Cuando ninguno de los dos es
conocido de antemano, el problema se formula como localización y mapeado
simultáneos, o \gls{slam}: la construcción incremental de un mapa de un entorno
desconocido a la vez que se estima la trayectoria del robot
\parencite{durrantwhyte2006slam}. Se trata de un problema fundamental en
robótica móvil porque ambas estimaciones son interdependientes, ya que un buen
mapa requiere conocer la posición y una buena posición requiere conocer el mapa,
y porque los errores de odometría se acumulan a lo largo del recorrido y
producen deriva si no se corrigen.

Una solución de \gls{slam} suele descomponerse en dos partes. El componente
frontal procesa las medidas de los sensores y resuelve la asociación de datos,
esto es, reconocer qué observaciones corresponden a un mismo elemento del
entorno; el componente posterior estima la trayectoria y el mapa más coherentes
con esas observaciones. La detección de cierres de bucle, que identifica lugares
ya visitados, es determinante para corregir la deriva acumulada. Históricamente,
las primeras formulaciones se apoyaban en filtros recursivos, como el filtro de
Kalman extendido, mientras que los enfoques predominantes en la actualidad
plantean el problema como la optimización de un grafo de poses, más preciso y
escalable para trayectorias largas.

Según el sensor empleado, se distingue entre el \gls{slam} visual, basado en
cámaras, y el \gls{slam} láser, basado en un sensor \gls{lidar} que mide
distancias mediante haces de luz. En el ámbito de \gls{ros}, una de las
soluciones más extendidas para interiores es \textit{slam\_toolbox}, que
implementa \gls{slam} láser basado en grafos de poses y permite tanto el mapeado
en línea como la posterior localización sobre el mapa generado, además de guardar
y recargar mapas previamente construidos \parencite{macenski2021slamtoolbox}. El
presente trabajo emplea este enfoque para que el robot construya el mapa de la
estancia durante una fase inicial de exploración.

%% =============================================================================
\section{Navegación autónoma y planificación de rutas de cobertura}
\label{sec:navegacion}

Una vez disponible el mapa, la navegación autónoma consiste en planificar y
ejecutar trayectorias seguras hacia los objetivos deseados evitando obstáculos.
En \gls{ros}, esta funcionalidad la proporciona \gls{nav2}, una pila de
navegación que integra un planificador global, un controlador local y un conjunto
de comportamientos de recuperación, coordinados mediante árboles de comportamiento
(\textit{behavior trees}) \parencite{macenski2020nav2}. El planificador global
calcula una ruta hasta el objetivo sobre una representación del entorno conocida
como mapa de costes (\textit{costmap}), que combina los obstáculos del mapa con
una zona de inflado de seguridad alrededor de ellos; el controlador local genera
las órdenes de velocidad que siguen esa ruta mientras esquiva obstáculos
imprevistos; y los comportamientos de recuperación intervienen cuando el robot
queda bloqueado.

El empleo de árboles de comportamiento, en lugar de una máquina de estados fija,
dota a la navegación de una lógica de decisión modular y fácil de modificar, en
la que pueden añadirse o reordenarse comportamientos sin reescribir el sistema.
\gls{nav2} expone, además, interfaces de alto nivel para enviar al robot a una
pose objetivo o a través de una secuencia ordenada de poses, una capacidad
especialmente adecuada para las tareas de patrulla, en las que el robot debe
recorrer una lista de puntos predefinidos.

El recorrido sistemático de un espacio se corresponde con el problema de
planificación de rutas de cobertura (\textit{coverage path planning}), que busca
una trayectoria que haga pasar al robot por todas las regiones accesibles del
entorno \parencite{galceran2013coverage}. Existen múltiples estrategias, desde las
descomposiciones celulares, que dividen el espacio en regiones simples recorridas
en barridos de ida y vuelta (\textit{boustrophedon}), hasta los métodos basados en
rejilla o en árboles de expansión, y la elección entre ellas responde al
compromiso entre completitud de la cobertura, longitud de la ruta y coste de
cálculo. En este trabajo se adopta una aproximación sencilla basada en una rejilla
de puntos sobre el espacio libre del mapa, recorrida en zig-zag, que resulta
suficiente para cubrir estancias cerradas de tamaño moderado.

%% =============================================================================
\section{Visión por computador y detección de objetos}
\label{sec:vision}

La detección de objetos consiste en localizar y clasificar las instancias de
interés presentes en una imagen, devolviendo para cada una un recuadro
delimitador (\textit{bounding box}) y una categoría con su grado de confianza. Los
métodos modernos se apoyan en redes neuronales convolucionales (\gls{cnn}),
capaces de aprender directamente de los datos las características visuales
relevantes. Tradicionalmente se distinguen dos familias: los detectores de dos
etapas, como \gls{cnnobj}, que primero proponen regiones candidatas y después las
clasifican, con alta precisión pero mayor coste computacional; y los detectores de
una sola etapa, que resuelven la localización y la clasificación en una única
pasada sobre la imagen.

La familia \gls{yolo} pertenece a esta segunda categoría y supuso un avance
decisivo al plantear la detección como un único problema de regresión sobre la
imagen completa, dividida en una rejilla en la que cada celda predice recuadros y
probabilidades de clase, lo que permitió alcanzar velocidades adecuadas para el
procesamiento en tiempo real \parencite{redmon2016yolo}. Un paso posterior de
supresión (\textit{non-maximum suppression}) descarta las
detecciones redundantes que se solapan sobre un mismo objeto. A lo largo de
sucesivas versiones, \gls{yolo} ha mejorado la precisión y la eficiencia y ha
reducido el tamaño de los modelos.

Para su despliegue en plataformas con recursos limitados resulta habitual emplear
las variantes más ligeras, que ofrecen un compromiso favorable entre precisión y
coste computacional, frente a las variantes mayores que requieren una \gls{gpu}
potente. En este trabajo se utiliza la versión \textit{nano} de \gls{yolo} en su
implementación de Ultralytics \parencite{jocher2023ultralytics}, ejecutada en el
propio robot para detectar la presencia de personas durante la patrulla. Este
procesamiento en el propio robot (\textit{edge computing}) evita transmitir vídeo de
forma continua a un servidor, reduce la carga de red y favorece la privacidad, al
analizar las imágenes en el propio dispositivo.

%% =============================================================================
\section{Arquitecturas orientadas a eventos}
\label{sec:arquitecturas-eventos}

En los sistemas distribuidos, una arquitectura orientada a eventos organiza la
comunicación en torno a la producción y el consumo de eventos a través de un
intermediario de mensajería (\textit{broker}), en lugar de mediante llamadas
directas entre componentes. Un evento es un registro inmutable de algo que ha
sucedido, por ejemplo una lectura de telemetría o la detección de una incidencia.
Este patrón desacopla a los productores de los consumidores: ambos desconocen su
existencia mutua y se comunican únicamente a través del bus, de modo que añadir un
nuevo consumidor no obliga a modificar al productor.

Una de las plataformas de referencia para este propósito es Apache Kafka. Utilizado por empresas como Netflix, Spotify o Uber, 
Kafka es un sistema de mensajería distribuido y persistente diseñado para procesar grandes
volúmenes de flujos de datos con baja latencia \parencite{kreps2011kafka}. Kafka
organiza los eventos en \textit{topics}, que a su vez se reparten en particiones
para escalar de forma horizontal, y conserva los mensajes en un registro
durante un periodo configurable. Esta persistencia permite que distintos
consumidores lean el mismo flujo a su propio ritmo, e incluso reprocesen eventos
pasados, una propiedad difícil de obtener con una comunicación síncrona y efímera.

Este enfoque resulta especialmente adecuado para un robot de patrulla. Al publicar
la telemetría y las incidencias como eventos, el robot solo transmite información
cuando se produce algo relevante y no necesita mantener conexiones síncronas con
el servidor ni esperar su respuesta, lo que reduce el consumo y mejora la
tolerancia a desconexiones temporales. El servidor, por su parte, consume estos
eventos a su propio ritmo y los distribuye al resto de la plataforma. Frente a un
modelo basado únicamente en peticiones \gls{rest}, en el que el cliente debe
preguntar de forma repetida por las novedades, el bus de eventos invierte el flujo
y reparte mejor la carga. Este desacoplamiento es uno de los pilares del diseño
propuesto.

%% =============================================================================
\section{Trabajos relacionados}
\label{sec:trabajos-relacionados}

La vigilancia mediante robots móviles autónomos cuenta con una trayectoria
consolidada en la literatura. Ya en propuestas tempranas se integraban el
mapeado, la localización y la navegación autónoma con tareas propias de la
vigilancia, como la detección de objetos abandonados o el seguimiento de
personas \parencite{dipaola2010surveillance}. Estos sistemas demostraron la
viabilidad de delegar la inspección de un entorno interior en una plataforma
autónoma y sentaron las bases sobre las que se apoyan los trabajos posteriores.

Las propuestas más recientes incorporan el aprendizaje profundo para la
percepción y se construyen sobre \gls{ros}. Existen robots de patrulla que
combinan \gls{slam} láser con cámaras RGB-D y emplean detectores de la familia
\gls{yolo} para generar mapas semánticos y evitar obstáculos durante el recorrido
\parencite{chiang2026patrol}. En la línea de la detección de intrusos, otras
propuestas despliegan flotas de robots coordinados bajo \gls{ros}, equipados con
sensores \gls{lidar} bidimensionales, que agregan sus detecciones en un nodo
central \parencite{islam2020intruder}. Estos trabajos confirman la idoneidad de
la combinación de \gls{ros}, \gls{slam} y detección basada en redes neuronales,
que también adopta el presente trabajo.

No obstante, la mayoría de estas soluciones concentran su esfuerzo en la propia
plataforma robótica, es decir, en la navegación y la percepción, y tratan de
forma secundaria, cuando no omiten, la integración con servicios web, la
persistencia y la gestión de las incidencias y la interacción con el usuario
final a través de una aplicación. Muchas son, además, sistemas cerrados o
difíciles de reproducir.

La propuesta de este trabajo se diferencia por abordar el sistema de forma
integral y con tecnologías abiertas: cubre desde la cartografía y la patrulla
autónoma hasta la notificación de incidencias y su consulta desde una aplicación
móvil, articulando todo ello mediante una arquitectura orientada a eventos. Esta
visión de extremo a extremo, más que la novedad de cada componente por separado,
constituye la aportación principal.

%% =============================================================================
\section{Síntesis y posicionamiento}
\label{sec:sintesis}

A lo largo del capítulo se han revisado tanto los fundamentos que sustentan el
trabajo, \gls{ros} como middleware, el \gls{slam} para el mapeado, \gls{nav2}
para la navegación, la detección de objetos con \gls{yolo} para la percepción y
las arquitecturas orientadas a eventos para la comunicación, como los trabajos
previos que abordan problemas afines. Cada una de esas tecnologías constituye una
pieza madura y consolidada en su ámbito, y los sistemas de vigilancia robótica
existentes confirman su aplicabilidad conjunta. Frente a esos trabajos, que
suelen detenerse en la plataforma robótica, el valor del presente reside en
integrar todas esas piezas, junto con el servidor y la aplicación de usuario, en
un único sistema funcional de extremo a extremo, cuyo diseño y construcción se
detallan en el Capítulo~\ref{cap:desarrollo}, tras formular los objetivos en el
Capítulo~\ref{cap:objetivos}.
